\documentclass[11pt]{article}

\usepackage[margin=1in]{geometry}
\usepackage{amsmath, amssymb, amsthm}
\usepackage{enumitem}

\pagestyle{empty}

\begin{document}

\begin{center}
  {\Large\bfseries Weekly Assignment: Direct Proof \&\\[2pt]
   Counterexample (Epp~4.1--4.4)}
\end{center}

\bigskip

\noindent
Complete each of the following problems.  Show all work and justify your answers.

\bigskip

%% ─────────────────────────────────────────────
\noindent\textbf{Problem 1} (Epp 4.1 \#2). \;
Assume that $c$ is a particular integer.  Justify your answers
by using the definitions of even and odd.

\begin{enumerate}[label=\textbf{\alph*.},itemsep=4pt]
  \item Is $-6c$ an even integer?
  \item Is $8c + 5$ an odd integer?
  \item Is $(c^2 + 1) - (c^2 - 1) - 2$ an even integer?
\end{enumerate}

\bigskip

%% ─────────────────────────────────────────────
\noindent\textbf{Problem 2} (Epp 4.1 \#9). \;
Prove the following statement:

\begin{quote}
  There is a real number $x$ such that $x > 1$ and $2^x > x^{10}$.
\end{quote}

\noindent
\emph{Hint}: If you're stuck, think about taking the
$\log$ of both sides.

\bigskip

%% ─────────────────────────────────────────────
\noindent\textbf{Problem 3} (Epp 4.1 \#13). \;
Consider the following statement:

\begin{quote}
  For every integer $n$, if $n$ is odd then
  $\displaystyle\frac{n^2 - 1}{2}$ is odd.
\end{quote}

\begin{enumerate}[label=\textbf{(\alph*)},itemsep=4pt]
  \item Write a negation for this statement.
  \item Use a counterexample to disprove the statement.
        Explain how the counterexample actually shows
        that the statement is false.
\end{enumerate}

\bigskip

%% ─────────────────────────────────────────────
\noindent\textbf{Problem 4} (Epp 4.2 \#5). \;
Prove the following statement directly from the definitions of
even and odd.  Follow the directions for writing proofs given in
Section~4.2.

\begin{quote}
  If $a$ and $b$ are any odd integers, then $a^2 + b^2$ is even.
\end{quote}

\bigskip

%% ─────────────────────────────────────────────
\noindent\textbf{Problem 5} (Epp 4.3 \#10). \;
Assume that $m$ and $n$ are both integers and that $n \ne 0$.
Explain why $\displaystyle\frac{5m^2 + 2n}{4n}$ must be a
rational number.

\bigskip

%% ─────────────────────────────────────────────
\noindent\textbf{Problem 6} (Epp 4.4 \#17). \;
Prove the following statement directly from the definition of
divisibility:

\begin{quote}
  For all integers $a$, $b$, $c$, and $d$, if $a \mid c$ and
  $b \mid d$ then $ab \mid cd$.
\end{quote}

\bigskip

%% ─────────────────────────────────────────────
\noindent\textbf{Problem 7} (Proof: Rotations as a Group). \;
Consider an object that can be rotated around its vertical axis,
either clockwise or counterclockwise.  We can think of all
possible rotations as a set~$G$, with the binary operation being
\emph{composition}: perform the first rotation, then perform the
second.  (One convenient representation is to identify each
rotation with its angle in degrees, so that $G$ is the set of all
real numbers modulo~$360$.)

Prove that this set of rotations forms a \textbf{group} under
composition.  Your proof should address each of the following
parts:

\begin{enumerate}[label=\textbf{(\alph*)},itemsep=6pt]
  \item \textbf{Identity.}  Define the identity element of this
        group and explain why composing any rotation with it
        leaves the rotation unchanged.

  \item \textbf{Inverses.}  Define the inverse of an arbitrary
        rotation and explain why composing a rotation with its
        inverse yields the identity.

  \item \textbf{Group axioms.}  Argue that the remaining group
        axioms hold:
        \begin{itemize}[itemsep=2pt]
          \item \emph{Closure}: composing any two rotations
                produces another rotation in~$G$.
          \item \emph{Associativity}: for any three rotations
                $r_1$, $r_2$, $r_3$,
                $(r_1 \circ r_2) \circ r_3
                 = r_1 \circ (r_2 \circ r_3)$.
        \end{itemize}
\end{enumerate}

\medskip\noindent
You may represent rotations by their degree measure (or
radians, or however you like).  We are looking for a solid
attempt at making the argument rather than any single particular
solution---this is practice at doing mathematics like a
mathematician!

\end{document}
